\subsection{I-33. Efficient prime cocharacters for two dense products, with a product-closure repair}
\label{I-33}

\paragraph{Scope and relation to known PIT.}
This construction targets efficient sampling rather than another
support bound.  Let $n\ge1$ and let $\delta\ge1$ be an integer;
the case $\delta=0$ consists only of constants.  Consider
\[
 f(\mathbf X)=\prod_{i=1}^n F_i(X_i)-\prod_{i=1}^n G_i(X_i),
 \qquad \deg F_i,\deg G_i\le\delta.
 \tag{I.62}
\]
It is a width-two diagonal ROABP, but its expanded support
may be exponential in $n$.  Polynomial PIT for this class
is already known, for example from I-B06.  Bisht, Gupta,
and Volkovich (2023), Theorem~1 (printed p.~9:4) and
Theorem~48 (printed p.~9:20), give a broader positive result
for sums of two products of read-once formulas.  We claim
no new PIT complexity for that established class and no
historical priority for the particular construction below.
The point is an explicit small family of compatible
integer cocharacters, a proof using divisors rather than
expanded monomials, and an exact test of its product closure.

Set
\[
 H=2\delta(n-1),\qquad M=H+2n+2,\qquad B=16M^2.
\]
Choose any $2n$ distinct integer primes
$p_1,\ldots,p_n,q_1,\ldots,q_n$ in $(H,B]$.
For $j=0,1,\ldots,n$, define positive integer weights
\[
 a_i^{(0)}=p_i,\qquad
 a_i^{(j)}=\begin{cases}q_j&i=j,\\p_i&i\ne j.\end{cases}
 \tag{I.63}
\]
These primes exist and can be found in polynomial time in
$n,\delta$ by an elementary sieve, as proved below.
Over \emph{every} field, every nonzero $f$ in (I.62) satisfies
\[
 f(T^{a_1^{(j)}},\ldots,T^{a_n^{(j)}})\ne0
 \quad\hbox{for at least one }0\le j\le n.
 \tag{I.64}
\]
No prime needs to be excluded when it equals the characteristic.
Each substituted polynomial has degree at most $n\delta B$.
Hence evaluating all $n+1$ curves on any $n\delta B+1$
distinct scalars gives a hitting set of size at most
\[
 S=(n+1)(n\delta B+1),
 \tag{I.65}
\]
with polynomial integer weights and polynomial query count.
As elsewhere, a working extension is required when the
coefficient field has too few scalar values.

\paragraph{1. The divisor separation lemma.}
Let $R_i(X)=F_i(X)/G_i(X)$ be nonzero rational functions,
with numerator and denominator degrees at most $\delta$,
and let $p_i$ be distinct primes greater than $H$.
If
\[
 \prod_i R_i(T^{p_i})\in\overline F^{\times},
 \tag{I.66}
\]
then each $R_i(X)$ is $c_iX^{b_i}$ for some nonzero
constant $c_i$ and integer $|b_i|\le\delta$.
Work over an algebraic closure and cancel numerator and
denominator common factors first.

If the characteristic equals one of the primes, say $p_k$,
then a nonzero root or pole of $R_k(T^{p_k})$ has nonzero
integer multiplicity of absolute value at least $p_k$.
At the same point, each other $R_i(T^{p_i})$ has multiplicity
of absolute value at most $\delta$, since the other power
maps are unramified away from zero.  Their combined
multiplicity has absolute value at most $(n-1)\delta<p_k$.
It therefore cannot cancel the first multiplicity, contrary
to (I.66).  Thus $R_k$ has no nonzero root or pole and is
a constant times a monomial.  There is at most one such
characteristic prime.

For any remaining $i$, a nonzero root or pole of $R_i$
pulls back to an orbit $O=\alpha\mu_{p_i}$ of $p_i$
distinct points.  Its multiplicity is nonzero at every
point of $O$.  For $j\ne i$ with $p_j$ different from the
characteristic, the nonzero divisor of $R_j(T^{p_j})$
is supported on at most $2\delta$ cosets of $\mu_{p_j}$.
Each such coset intersects $O$ in at most one point, since
$\gcd(p_i,p_j)=1$.  A characteristic-prime factor, if
present, was already shown to have no nonzero divisor.
Consequently all other factors together can cancel the
divisor on at most $2\delta(n-1)=H<p_i$ points of $O$.
This contradicts (I.66).

Every $R_i$ therefore has divisor supported only at zero
and infinity.  Factoring its numerator and denominator
over the algebraic closure gives $R_i(X)=c_iX^{b_i}$.
The degree bounds give $|b_i|\le\delta$.  This proves
the lemma, including the inseparable characteristic case.

\paragraph{2. The family separates the remaining monomial relation.}
If one product in (I.62) is zero and the other is nonzero,
every positive-weight curve keeps the nonzero product
nonzero, so (I.64) is immediate.  Otherwise both products
are nonzero and we may divide by the second.
If the base curve $j=0$ kills $f$, the divisor lemma gives
$R_i(X_i)=c_iX_i^{b_i}$, and the resulting Laurent identity
forces
\[
 \prod_i c_i=1,\qquad\sum_i b_ip_i=0.
\]
If every $b_i=0$, then $f=0$, contrary to the hypothesis.
Choose $j$ with $b_j\ne0$.  Changing only its prime changes
the exponent sum by the nonzero integer
$b_j(q_j-p_j)$.  Thus the $j$th curve does not kill $f$.
This is an equality of exponents of the formal indeterminate
$T$, so it is valid even when that integer is divisible by
the characteristic.  This proves (I.64).

The same argument shows that if the two component products
are linearly independent, at least one curve preserves
their two-dimensional coefficient span.  If their ratio is
not a monomial product, the base-curve ratio cannot be
constant by the divisor lemma.  If it is a nonconstant
monomial product, some weight vector has nonzero exponent
sum.  In either case the two univariate images are linearly
independent.  This is a concrete instance of the compression
condition I-14; no basis-isolating-weight claim is needed.

The weight family has full affine span: its differences
from the base vector are $(q_j-p_j)e_j$.  Encoding
$j\in\{0,\ldots,n\}$ by $\lceil\log_2(n+1)\rceil$ bits
gives an explicit nonlinear decoder using coordinate
indicator functions.  It is therefore outside the bounded
affine and fixed-degree decoder obstruction of I-15.
It is not obtained from Hoza's path-isolation theorem.

\paragraph{3. An elementary polynomial prime bound.}
For even $B$, the central binomial coefficient satisfies
\[
 \binom{B}{B/2}\ge\frac{2^B}{B+1}.
\]
For each prime $p\le B$, its valuation in this binomial
coefficient is at most $\lfloor\log_p B\rfloor$: each
term in the factorial-valuation difference is zero or one.
Hence its prime-power contribution is at most $B$, and
\[
 B^{\pi(B)}\ge\frac{2^B}{B+1},\qquad
 \pi(B)\ge\frac{B\log2-\log(B+1)}{\log B}.
\]
For $B=16M^2$ and $M\ge2$, the elementary inequalities
$\log(B+1)\le(B/2)\log2$ and $\log B\le4M\log2$
give $\pi(B)\ge2M$.  At most $H$ primes are at most $H$,
so more than $2n$ primes lie in $(H,B]$.
Sieve this polynomially bounded interval and take its first
$2n$ primes.  This avoids using an unproved prime-distribution
hypothesis or a bit-length-only prime bound.

\paragraph{4. The raw curve family is not closed under products.}
The preceding success cannot be extended to general width
two just by multiplying separate successful factors.
For the explicit construction with $n=64$ and $\delta=1$,
one has $B=16\cdot256^2$.  Split the variables into two
blocks of size $32$.  Within either block the $2^{32}$
subset sums of the base primes take at most $32B+1$
integer values, and $2^{32}>32B+1$.  Two different subsets
therefore have the same sum.  After canceling their common
indices, this gives a nonzero multilinear binomial on that
block which vanishes along the base-prime curve.
Let the two disjoint-block binomials be $A$ and $C$.
Their product $AC$ is nonzero and multilinear.

Changing one coordinate's prime affects at most one of
$A,C$; the other remains zero along that curve.  Thus
$AC$ vanishes along all $n+1$ curves in (I.63).
It is nevertheless a width-two ROABP in a suitable order:
each binomial is the difference of two monomial products,
computed by diagonal $2\times2$ factors, and their two
scalar computations concatenate through a constant
rank-one bridge.  Absorbing that bridge into an adjacent
univariate layer gives an ordinary width-two ROABP.
This is a counterexample to the specified raw family,
not to the known width-two hitting-set theorem I-B06.

\paragraph{5. A concrete product-closure repair.}
Let $\mathcal H$ be the list of at most $S$ points from
(I.65).  Interpolate a polynomial curve
$\Gamma(Z)\in E[Z]^n$ of coordinate degree at most $S-1$
through that list at $S$ distinct nodes, repeating points
if necessary.  Every nonzero factor of the form (I.62)
stays nonzero on $\Gamma$, because some listed point hits it.
Therefore any nonzero product of such factors also stays
nonzero, even when different factors share variables.
For a product of total degree at most $D\ge1$, its
composition has degree at most $D(S-1)$ and is hit on any
$D(S-1)+1$ distinct scalars.  A working field with at least
that many elements supplies both the interpolation nodes
and this final grid.

This is the scalar product-closure mechanism of AGKS,
Lemma~33, recorded in I-B06, applied to the new explicit
prime-curve family.  It repairs the precise four-term
counterexample above and gives a fully polynomial sampler
for products of the dense two-product factors.  It does
not handle arbitrary sums of three or more products,
general noncommuting matrix blocks, or their coefficient
matroids.  Those cancellation patterns remain outside
the divisor-separation argument.
\hfill$\square$
